%  The 24C serial-EEPROM family: one generic driver over the whole catalogue.
\chapter[Serial EEPROM: the 24C Family]%
        {Serial EEPROM: the 24C Family --- One Driver, a Whole Catalogue}
\label{ch:eeprom}

An SD card holds gigabytes and a NOR flash holds megabytes. Sometimes you want
kilobytes: a serial number, a calibration table, a boot counter. That is what a
24C EEPROM is for. It is byte-addressable, needs no erase before a write, and
survives far more write cycles than flash. It costs a few cents and speaks I2C.

The family is large. ST sells the M24Cxx line, Microchip the 24LCxx and 24AAxx,
Atmel the AT24Cxx, onsemi the CAT24Cxx. Every one of them speaks the same
protocol. Three details make that protocol a trap for the unwary, and
\texttt{ESP32S3.EEPROM\_24C} handles all three.\index{EEPROM}\index{24C family}

\textbf{A page write wraps.}\index{page write} The part programs 32 bytes at a
time (fewer on small parts, more on large ones). A write that runs past a page
boundary does not advance into the next page. It wraps to the start of the
current one and overwrites what you just sent. Nothing complains. \texttt{Write}
splits its payload on page boundaries so this cannot happen.

\textbf{A write takes 5\,ms, and the part goes deaf.}\index{ACK polling} Once a
page write lands, the chip starts an internal program cycle. Until it finishes, it
does not acknowledge its own address. The fix is \emph{ACK polling}: keep probing
the address until it answers. \texttt{Write} polls rather than sleeping the worst
case, so a fast part does not pay for a slow one.

\textbf{A random read needs a repeated START.} You write the memory address, then
turn the bus around without a STOP. A STOP would end the command. Every datasheet
in the family draws it this way, and the driver issues it as one
\texttt{ESP32S3.I2C.Write\_Read} (Section~\ref{driver:end-opcode}).

\section{The API}

What remains, once those three traps are handled inside the driver, is a plain,
boring API. \texttt{Read} and \texttt{Write} take an arbitrary
\texttt{Byte\_Array} at an arbitrary address; \texttt{Read\_Byte} /
\texttt{Write\_Byte} are the single-cell shorthands. Every operation returns a
\texttt{Status} --- \texttt{OK}, \texttt{Bus\_Error} (no ACK: absent,
mis-strapped, or write-protected), or \texttt{Write\_Timeout} (the program cycle
never finished). \texttt{Is\_Present} is an address-only probe.

\begin{lstlisting}
with ESP32S3.EEPROM_24C.M24C64;
...
   package Rom_Part renames ESP32S3.EEPROM_24C.M24C64;
   Rom : Rom_Part.Device;
   St  : Rom_Part.Status;
   Buf : ESP32S3.I2C.Byte_Array (0 .. 63);
begin
   Rom_Part.Setup (Rom, Sda => 41, Scl => 40);   --  A0/A1/A2 strapped low
   if Rom_Part.Is_Present (Rom) then
      Rom_Part.Write (Rom, 16#0100#, Buf, St);    --  splits on page boundaries
      Rom_Part.Read  (Rom, 16#0100#, Buf, St);    --  one transaction, any length
   end if;
\end{lstlisting}

A \texttt{Read} holds the I2C host for the whole transfer; a \texttt{Write} holds
it across every page \emph{and} its program cycle, so a multi-page \texttt{Write}
is atomic with respect to other tasks on the bus. On a \texttt{Bus\_Error} or
\texttt{Write\_Timeout} mid-write, the pages before the failure are already
committed --- the array is not rolled back.

The write-protect pin --- \texttt{WC} on ST parts, \texttt{WP} on everyone
else --- is \emph{wiring, not software}: held HIGH it inhibits every write in the
family, and there is no register for it. Tie it low to allow writes. Because the
driver rides the controlled I2C \texttt{Session} (with finalization), it is an
embedded-/full-profile driver.

\section{A generic over the family}\index{generic package}\index{catalogue}

Since all the parts share a protocol, they should share a driver. What differs is
only geometry: how big the array is, how big a page is, and whether the word
address takes one byte or two. So the driver is a generic, and a part is a value:

\begin{lstlisting}
type Geometry is record
   Capacity_Bytes     : Positive;
   Page_Bytes         : Positive;
   Word_Address_Bytes : Positive;   --  1 up to 16 Kbit, 2 from 32 Kbit up
   Max_Read_Span      : Natural := 0;
   Base_Slave_Address : ESP32S3.I2C.Slave_Address := 16#50#;
   Tested             : Verification := Untested;
end record;

M24C64_Part : constant Geometry :=
  (Capacity_Bytes => 8_192, Page_Bytes => 32, Word_Address_Bytes => 2,
   Tested => Verified, others => <>);
\end{lstlisting}

\texttt{ESP32S3.EEPROM\_24C} is the catalogue: one package holding a
\texttt{Geometry} for every part. \texttt{ESP32S3.EEPROM\_24C.Driver} is the
generic. Each part is then a child unit one line long:

\begin{lstlisting}
package ESP32S3.EEPROM_24C.M24C16 is
  new ESP32S3.EEPROM_24C.Driver (M24C16_Part);
\end{lstlisting}

Child units, not packages nested in the catalogue. The difference is real: a
nested instantiation would compile every part into every program that touched the
catalogue, about 3\,kB each. As children, \texttt{with}-ing one part costs one
part. The catalogue ships the ST line M24C01 through M24M02, the Atmel AT24C01 /
AT24C02, and Microchip's 24LC1026 --- from 128 bytes to 256\,KiB.

\section{The fourth field that is not there}\index{chip-enable strap}

A 24C part carries three chip-enable pins, E2, E1 and E0, strapped high or low on
the board. They form the low three bits of the slave address, so eight parts can
share a bus: \texttt{16\#50\#} through \texttt{16\#57\#}.

Except when they do not. A part whose array outruns its word address must put the
surplus address bits somewhere. It puts them in the low bits of its own
device-select byte, and each one eats a chip-enable pin, starting at E0. So a
24C04 folds A8 and keeps two straps. A 24C16 folds A10, A9 and A8, keeps none, and
only one 24C16 can ever sit on a bus.

That could have been a fourth field in \texttt{Geometry}, and a fourth chance for
someone to get it wrong. It is derived instead:

\begin{lstlisting}
Word_Span   : constant Positive := 2 ** (8 * Address_Bytes);   --  256 or 65536
Blocks      : constant Positive := (Capacity + Word_Span - 1) / Word_Span;
Max_Devices : constant Positive := 8 / Blocks;

--  and the slave address for the byte at Location folds in the high bits:
function Slave_Of (Dev : Device; Location : Memory_Address) return Slave_Address
is (Base_Address + Dev.Strap + Location / Word_Span);
\end{lstlisting}

The folded bits and the straps never collide, because a strap the part has
swallowed is not a pin any more. \texttt{Setup} says so in a precondition, and the
precondition is per instance:

\begin{lstlisting}
procedure Setup (Dev : out Device; ...; A0, A1, A2 : Pin_State := Low; ...)
with Pre => (if Blocks >= 2 then A0 = Low)
            and then (if Blocks >= 4 then A1 = Low)
            and then (if Blocks >= 8 then A2 = Low);
\end{lstlisting}

\noindent \texttt{M24C64.Setup (Rom, A0 => High)} is fine and gives you
\texttt{16\#51\#}. \texttt{M24C16.Setup (Rom, A0 => High)} raises, because that pin
does not exist on that part. One generic, two contracts, no runtime cost.

\section{Marking what has not been tested}

Only the M24C64 has been run against real silicon here. The rest are transcribed
from datasheets. The protocol is shared, so they are very likely right, but
``very likely'' is not the same as ``seen on a scope.''

The catalogue says so. Each \texttt{Geometry} carries \texttt{Tested}, each
instance re-exports \texttt{Hardware\_Verified : constant Boolean}, and each
part's spec opens with a \texttt{STATUS: UNTESTED} banner. This is cheap honesty,
and it beats a comment nobody updates. Flip a part to \texttt{Verified} (and say
on what board) once you have run it.

Two traps are worth naming, because both corrupt data quietly. ST's 1\,Kbit and
2\,Kbit parts use a 16-byte page; Atmel's and Microchip's use 8. Guess wrong and
the part wraps mid-write. Hence separate \texttt{AT24C01} and \texttt{AT24C02}
entries. And Microchip's 24LC1025 puts its block-select bit in the \emph{high}
position of the device-select byte, not the low one, so this driver's addressing
model cannot express it. It is deliberately absent, rather than present and wrong.
Its sibling the 24LC1026 folds normally, and is there --- with a
\texttt{Max\_Read\_Span} so its sequential read is split at the 512-Kbit block
boundary it may not cross.

\noindent The on-target example \texttt{esp32s3\_m24c64} probes the strapped
address, keeps a boot counter across resets, and writes a pattern across a page
boundary to prove the split.
